\chapter{Les alcanes}

\textit{Les alcanes constituent la famille la plus simple des hydrocarbures saturés. Présents dans le pétrole et le gaz naturel, ils sont à la base de nombreux combustibles et matières premières de l'industrie pétrochimique. Leur étude permet de comprendre les concepts fondamentaux de la chimie organique : nomenclature, isomérie, réactions radicalaires et transformations industrielles.}

\section{Définition et série homologue}

\begin{definition}[Alcane]
Un alcane est un hydrocarbure saturé, c'est-à-dire un composé organique ne contenant que des atomes de carbone et d'hydrogène, dans lequel toutes les liaisons entre atomes de carbone sont des liaisons simples. La formule générale des alcanes non cycliques (alcanes linéaires ou ramifiés) est :
\[
\boxed{\ce{C_nH_{2n+2}}}
\]
où \(n\) est un nombre entier naturel non nul (\(n \ge 1\)).
\end{definition}

\begin{propriete}[Série homologue]
Les alcanes forment une série homologue : chaque terme diffère du précédent par un groupe \ce{CH2}. Les premiers termes sont :
\begin{itemize}
    \item \ce{CH4} : méthane (\(n=1\))
    \item \ce{C2H6} : éthane (\(n=2\))
    \item \ce{C3H8} : propane (\(n=3\))
    \item \ce{C4H10} : butane (\(n=4\))
    \item \ce{C5H12} : pentane (\(n=5\))
\end{itemize}
\end{propriete}

\begin{exemple}
Le décane, de formule brute \ce{C10H22}, est le dixième terme de la série (\(n=10\)).
\end{exemple}

\section{Nomenclature des alcanes ramifiés}

La nomenclature systématique (IUPAC) permet de nommer sans ambiguïté tout alcane ramifié.

\begin{regle}[Règles de nomenclature]
\begin{enumerate}
    \item \textbf{Chaîne principale} : choisir la chaîne carbonée la plus longue (nombre maximal d'atomes de C). Le nom de base est celui de l'alcane linéaire correspondant.
    \item \textbf{Numérotation} : numéroter la chaîne principale de façon à donner les indices les plus petits possibles aux ramifications (substituants).
    \item \textbf{Substituants} : les groupes alkyles (provenant d'un alcane par perte d'un H) sont nommés en remplaçant le suffixe \textit{-ane} par \textit{-yle}. Exemples : \ce{CH3-} (méthyle), \ce{C2H5-} (éthyle), \ce{CH3CH2CH2-} (propyle).
    \item \textbf{Ordre alphabétique} : les substituants sont cités dans l'ordre alphabétique (sans tenir compte des préfixes multiplicateurs di-, tri-, etc.).
    \item \textbf{Préfixes multiplicateurs} : si un même substituant apparaît plusieurs fois, on utilise \textit{di-}, \textit{tri-}, \textit{tétra-}, etc.
    \item \textbf{Séparation} : les nombres sont séparés par des virgules, et les nombres des lettres par des tirets.
\end{enumerate}
\end{regle}

\begin{exemple}[Nommer un alcane ramifié]
\begin{center}
\chemfig{CH_3-CH(CH_3)-CH_2-CH_3}
\end{center}
\begin{itemize}
    \item Chaîne la plus longue : 4 carbones (butane).
    \item Ramification : groupe méthyle sur le carbone \textnumero 2.
    \item Nom : \textbf{2-méthylbutane}.
\end{itemize}
\end{exemple}

\begin{exemple}[Cas de plusieurs ramifications]
\begin{center}
\chemfig{CH_3-CH(CH_3)-C(CH_3)(CH_3)-CH_2-CH_3}
\end{center}
\begin{itemize}
    \item Chaîne la plus longue : 5 carbones (pentane).
    \item Ramifications : deux groupes méthyle en position 2, deux groupes méthyle en position 3.
    \item Nom : \textbf{2,2,3,3-tétraméthylpentane}.
\end{itemize}
\end{exemple}

\section{Isomérie de chaîne}

\begin{definition}[Isomères de chaîne]
Les isomères de chaîne sont des composés ayant la même formule brute mais des enchaînements différents des atomes de carbone (chaîne linéaire ou ramifiée).
\end{definition}

\begin{exemple}[Isomères du butane \ce{C4H10}]
\begin{itemize}
    \item \textbf{Butane} (chaîne linéaire) : \chemfig{CH_3-CH_2-CH_2-CH_3}
    \item \textbf{2-méthylpropane} (chaîne ramifiée) : \chemfig{CH_3-CH(CH_3)-CH_3}
\end{itemize}
\end{exemple}

\begin{propriete}
Le nombre d'isomères de chaîne augmente rapidement avec le nombre d'atomes de carbone :
\begin{itemize}
    \item \ce{C4H10} : 2 isomères
    \item \ce{C5H12} : 3 isomères
    \item \ce{C6H14} : 5 isomères
    \item \ce{C7H16} : 9 isomères
    \item \ce{C10H22} : 75 isomères
\end{itemize}
\end{propriete}

\section{Propriétés physiques}

\begin{propriete}[Propriétés physiques des alcanes]
\begin{itemize}
    \item \textbf{État physique} : à température ambiante, les alcanes de \(n=1\) à \(n=4\) sont gazeux, de \(n=5\) à \(n=17\) sont liquides, au-delà de \(n=18\) sont solides.
    \item \textbf{Point d'ébullition} : augmente régulièrement avec la masse molaire (plus la chaîne est longue, plus les forces de van der Waals sont importantes). Pour un même nombre de carbones, l'isomère le plus ramifié a le point d'ébullition le plus bas.
    \item \textbf{Densité} : inférieure à celle de l'eau (les alcanes flottent).
    \item \textbf{Solubilité} : insolubles dans l'eau (molécules apolaires), solubles dans les solvants organiques apolaires (éther, benzène, etc.).
    \item \textbf{Réactivité} : les alcanes sont peu réactifs (paraffines = « peu d'affinité »). Ils ne réagissent pas avec les acides, les bases, les oxydants doux.
\end{itemize}
\end{propriete}

\begin{experience}[Solubilité des alcanes]
Dans un tube à essai contenant de l'eau distillée, ajouter quelques gouttes d'hexane. Agiter. On observe deux phases : l'hexane (moins dense) surnage. L'hexane est insoluble dans l'eau.
\end{experience}

\section{Propriétés chimiques}

\subsection{Combustion}

\begin{definition}[Combustion]
La combustion est une réaction d'oxydation rapide d'un combustible (ici un alcane) par le dioxygène de l'air, avec dégagement de chaleur et de lumière.
\end{definition}

\begin{propriete}[Combustion complète]
En présence d'un excès de dioxygène, la combustion complète d'un alcane produit du dioxyde de carbone et de l'eau :
\[
\ce{C_nH_{2n+2} + \frac{3n+1}{2} O2 -> n CO2 + (n+1) H2O}
\]
\end{propriete}

\begin{exemple}[Combustion du méthane]
\[
\ce{CH4 + 2 O2 -> CO2 + 2 H2O}
\]
\end{exemple}

\begin{propriete}[Combustion incomplète]
En présence d'un défaut de dioxygène, la combustion est incomplète. On peut obtenir :
\begin{itemize}
    \item du monoxyde de carbone \ce{CO} (gaz toxique) :
    \[
    \ce{2 CH4 + 3 O2 -> 2 CO + 4 H2O}
    \]
    \item du carbone (suie) :
    \[
    \ce{CH4 + O2 -> C + 2 H2O}
    \]
\end{itemize}
\end{propriete}

\begin{remarque}
La combustion incomplète est dangereuse car le monoxyde de carbone \ce{CO} est inodore, incolore et très toxique (se fixe sur l'hémoglobine).
\end{remarque}

\subsection{Halogénation radicalaire}

\begin{definition}[Halogénation]
L'halogénation est la substitution d'un ou plusieurs atomes d'hydrogène d'un alcane par des atomes d'halogène (principalement le chlore \ce{Cl2} ou le brome \ce{Br2}), sous l'action de la lumière (UV) ou de la chaleur.
\end{definition}

\begin{experience}[Chloration du méthane]
Dans un flacon transparent contenant un mélange de méthane et de dichlore, exposer à la lumière du jour ou à une lampe UV. On observe une décoloration progressive du mélange (le dichlore vert-jaune disparaît) et formation de fumées blanches de chlorure d'hydrogène.
\end{experience}

\begin{propriete}[Équation bilan]
\[
\ce{CH4 + Cl2 ->[h\nu] CH3Cl + HCl}
\]
La réaction ne s'arrête pas au monochlorométhane : on obtient un mélange de produits chlorés (\ce{CH3Cl}, \ce{CH2Cl2}, \ce{CHCl3}, \ce{CCl4}).
\end{propriete}

\subsubsection{Mécanisme radicalaire en chaîne}

Le mécanisme de l'halogénation radicalaire comporte trois étapes :

\begin{enumerate}
    \item \textbf{Initiation} : rupture homolytique de la liaison \ce{Cl-Cl} sous l'effet de la lumière UV.
    \[
    \ce{Cl2 ->[h\nu] 2 Cl$^{\bullet}$}
    \]
    (Les points représentent des électrons célibataires : radicaux libres.)
    
    \item \textbf{Propagation} : deux réactions se répètent en chaîne.
    \begin{itemize}
        \item Arrachement d'un atome d'hydrogène :
        \[
        \ce{Cl$^{\bullet}$ + CH4 -> HCl + CH3$^{\bullet}$}
        \]
        \item Réaction du radical alkyle avec une molécule de dichlore :
        \[
        \ce{CH3$^{\bullet}$ + Cl2 -> CH3Cl + Cl$^{\bullet}$}
        \]
    \end{itemize}
    Le radical chlore \ce{Cl$^{\bullet}$} est régénéré, permettant la poursuite de la chaîne.
    
    \item \textbf{Terminaison} : recombination de deux radicaux pour former une molécule stable.
    \begin{itemize}
        \item \ce{Cl$^{\bullet}$ + Cl$^{\bullet}$ -> Cl2}
        \item \ce{CH3$^{\bullet}$ + Cl$^{\bullet}$ -> CH3Cl}
        \item \ce{CH3$^{\bullet}$ + CH3$^{\bullet}$ -> C2H6}
    \end{itemize}
\end{enumerate}

\begin{aretenir}
Le mécanisme radicalaire de l'halogénation des alcanes comporte trois phases : \textbf{initiation} (création de radicaux), \textbf{propagation} (réactions en chaîne), \textbf{terminaison} (disparition des radicaux).
\end{aretenir}

\subsection{Craquage et reformage}

\begin{definition}[Craquage]
Le craquage est un procédé pétrochimique qui consiste à casser les longues chaînes carbonées des alcanes lourds (provenant du pétrole) en molécules plus légères et plus utiles (essence, gaz).
\end{definition}

\begin{propriete}[Craquage thermique]
Sous l'effet de la chaleur (température élevée, \SI{500}{\celsius} à \SI{700}{\celsius}) et éventuellement d'un catalyseur, une molécule d'alcane se fragmente :
\[
\ce{C16H34 -> C8H18 + C8H16}
\]
(hexadécane $\rightarrow$ octane + octène)
\end{propriete}

\begin{definition}[Reformage]
Le reformage (ou reforming) est un procédé qui transforme les alcanes linéaires en alcanes ramifiés ou en hydrocarbures aromatiques, améliorant ainsi l'indice d'octane de l'essence.
\end{definition}

\begin{exemple}[Reformage]
\[
\ce{CH3(CH2)6CH3 ->[catalyseur][\Delta] C8H18\text{ (ramifié)}}
\]
ou
\[
\ce{C6H14 ->[catalyseur][\Delta] C6H6 + 4 H2}
\]
(hexane $\rightarrow$ benzène + dihydrogène)
\end{exemple}

\section{Le pétrole et le gaz naturel}

\begin{definition}[Pétrole brut]
Le pétrole brut est un mélange complexe d'hydrocarbures (principalement des alcanes, cycloalcanes et composés aromatiques) provenant de la décomposition de matière organique fossile.
\end{definition}

\begin{definition}[Gaz naturel]
Le gaz naturel est un mélange gazeux composé essentiellement de méthane (\ce{CH4}, 70-90\%), d'éthane, de propane et de butane.
\end{definition}

\begin{propriete}[Distillation fractionnée du pétrole]
Le pétrole brut est séparé en fractions par distillation fractionnée (raffinage). Chaque fraction correspond à une gamme de températures d'ébullition :
\begin{itemize}
    \item \textbf{Gaz} (\SI{<20}{\celsius}) : \ce{C1} à \ce{C4} (méthane, éthane, propane, butane)
    \item \textbf{Essence} (\SIrange{20}{200}{\celsius}) : \ce{C5} à \ce{C10}
    \item \textbf{Kérosène} (\SIrange{150}{250}{\celsius}) : \ce{C10} à \ce{C16}
    \item \textbf{Gazole} (\SIrange{250}{350}{\celsius}) : \ce{C14} à \ce{C20}
    \item \textbf{Résidus} (\SI{>350}{\celsius}) : \ce{C20} et plus (fioul lourd, bitume)
\end{itemize}
\end{propriete}

\begin{remarque}
Le craquage et le reformage permettent de valoriser les fractions lourdes en les transformant en hydrocarbures plus légers (essence) ou en améliorant leur qualité.
\end{remarque}

\section{L'essentiel à retenir}

\begin{synthese}
\subsection*{Les alcanes : fiche de révision}
\begin{itemize}
    \item \textbf{Formule générale} : \ce{C_nH_{2n+2}} (alcanes non cycliques).
    \item \textbf{Nomenclature} : chaîne la plus longue, numérotation pour indices minimaux, substituants en ordre alphabétique.
    \item \textbf{Isomérie} : isomères de chaîne (même formule brute, enchaînement différent).
    \item \textbf{Propriétés physiques} : apolaires, insolubles dans l'eau, \(T_{éb}\) croît avec \(n\).
    \item \textbf{Combustion} : complète (\ce{CO2 + H2O}) ou incomplète (\ce{CO} ou \ce{C}).
    \item \textbf{Halogénation radicalaire} : mécanisme en 3 étapes (initiation, propagation, terminaison).
    \item \textbf{Craquage} : fragmentation des longues chaînes en molécules plus légères.
    \item \textbf{Reformage} : transformation des alcanes linéaires en ramifiés ou aromatiques.
    \item \textbf{Pétrole} : mélange d'hydrocarbures séparé par distillation fractionnée.
\end{itemize}
\end{synthese}

\section{Méthodes \& savoir-faire}

\methode{Écrire la formule brute d'un alcane à partir de son nom}
Pour un alcane linéaire, le nom indique le nombre \(n\) de carbones (méth- = 1, éth- = 2, prop- = 3, but- = 4, pent- = 5, hex- = 6, hept- = 7, oct- = 8, non- = 9, déc- = 10). Appliquer la formule \ce{C_nH_{2n+2}}.

\begin{exemple}
Nom : heptane. \(n=7\), donc formule brute : \ce{C7H16}.
\end{exemple}

\methode{Nommer un alcane ramifié}
\begin{enumerate}
    \item Identifier la chaîne carbonée la plus longue.
    \item Numéroter pour donner les indices les plus petits aux ramifications.
    \item Citer les substituants par ordre alphabétique avec leur position.
    \item Écrire le nom en un seul mot (chiffres et lettres séparés par des tirets).
\end{enumerate}

\begin{exemple}
\chemfig{CH_3-CH_2-CH(CH_3)-CH(CH_3)-CH_3}
\begin{itemize}
    \item Chaîne la plus longue : 5 carbones (pentane).
    \item Numérotation : de gauche à droite (ramifications en 2 et 3).
    \item Noms : 2-méthyl, 3-méthyl.
    \item Nom : \textbf{2,3-diméthylpentane}.
\end{itemize}
\end{exemple}

\methode{Équilibrer une équation de combustion complète}
Utiliser la formule générale : \ce{C_nH_{2n+2} + \frac{3n+1}{2} O2 -> n CO2 + (n+1) H2O}. Multiplier par 2 si nécessaire pour éviter les fractions.

\begin{exemple}
Combustion du propane (\ce{C3H8}) : \(n=3\), donc \ce{C3H8 + 5 O2 -> 3 CO2 + 4 H2O}.
\end{exemple}

\methode{Écrire le mécanisme de l'halogénation radicalaire}
\begin{enumerate}
    \item \textbf{Initiation} : \ce{Cl2 ->[h\nu] 2 Cl$^{\bullet}$}.
    \item \textbf{Propagation} : \ce{Cl$^{\bullet}$ + CH4 -> HCl + CH3$^{\bullet}$} puis \ce{CH3$^{\bullet}$ + Cl2 -> CH3Cl + Cl$^{\bullet}$}.
    \item \textbf{Terminaison} : recombination de deux radicaux (ex. \ce{CH3$^{\bullet}$ + Cl$^{\bullet}$ -> CH3Cl}).
\end{enumerate}

\begin{exemple}
Pour la chloration de l'éthane, le radical éthyle \ce{C2H5$^{\bullet}$} se forme en propagation.
\end{exemple}

\methode{Distinguer combustion complète et incomplète}
\begin{itemize}
    \item Complète : flamme bleue, produits \ce{CO2} et \ce{H2O}.
    \item Incomplète : flamme jaune/suie, présence de \ce{CO} ou \ce{C}.
\end{itemize}

\begin{exemple}
Un bec Bunsen mal réglé donne une flamme jaune (dépôt de suie) : combustion incomplète du méthane.
\end{exemple}

\methode{Calculer le nombre d'isomères de chaîne}
Pour les petits alcanes, les dessiner systématiquement en commençant par la chaîne linéaire, puis en raccourcissant la chaîne principale et en plaçant les ramifications.

\begin{exemple}
Pour \ce{C5H12} (pentane) : 3 isomères (pentane, 2-méthylbutane, 2,2-diméthylpropane).
\end{exemple}

\methode{Interpréter la distillation fractionnée du pétrole}
Associer chaque fraction à sa plage de température et à ses utilisations.

\begin{exemple}
La fraction essence (\SIrange{20}{200}{\celsius}) contient des alcanes en \ce{C5-C10} utilisés comme carburant.
\end{exemple}

\methode{Prévoir les produits d'un craquage}
Le craquage d'un alcane lourd donne un alcane plus léger et un alcène.

\begin{exemple}
\ce{C12H26 -> C6H14 + C6H12} (dodécane $\rightarrow$ hexane + hexène).
\end{exemple}

\methode{Utiliser la nomenclature des groupes alkyles}
Reconnaître les groupes alkyles courants : méthyle (\ce{CH3-}), éthyle (\ce{C2H5-}), propyle (\ce{CH3CH2CH2-}), isopropyle (\ce{(CH3)2CH-}).

\begin{exemple}
Le 2-méthylpropane a un groupe méthyle sur le carbone 2 du propane.
\end{exemple}